\SourceLeaf{524}forma contentorum multitudo sit $=\frac12(p-1)p^{\pi-3}$; pro singulis autem fit $M=2p$. Summa terminorum in $\Pi$ hinc oriundorum fit
\[
\begin{aligned}
&=2p\Sigma r^{\mu pp+kp^3}=2p\Sigma r^{\mu pp}\pdot\Sigma r^{kp^3}\\
&=2p\Sigma r^{\mu pp}\pdot\frac{r^{p^\pi}-1}{r^{p^3}-1}=0
\end{aligned}
\]
siquidem $\pi>3$. Similiter classis tertia, quarta etc. exhibebitur per $\mu p^4+kp^5$, $\mu p^6+kp^7$ etc. ubi pro $k$ omnes integri a $0$ usque ad $p^{\pi-5}-1$, $p^{\pi-7}-1$ etc. accipi debent; pro his fit $M=2pp$, $M=2p^3$ etc. Et summa terminorum in $\Pi$ e classe tertia, quarta etc. ortorum evanescet, siquidem $\pi>5$, $\pi>7$ etc.. resp.
\par Hinc colligitur, pro casu ubi $\pi$ par est, in $\Pi$ eos tantummodo terminos remanere, qui residuo $0$ respondent, qui sunt $=1$; pro his vero fit $M=p^{\frac12\pi}$, ita ut summa omnium terminorum in $\Pi$ fiat $=p^{\frac12\pi}$.
\BookRule{.27\textwidth}
